% GENERATED FILE. Edit canonical lessons, then run npm run generate:books.
% canonical-source-hash: fnv1a64:8f97467d746dbfe5
% canonical-lessons: TE-C481-vadagali, TE-C481-endamavi, TE-C481-galidumaram, TE-C481-endipovu, TE-C481-bitalu-varu

\chapter{Heat and Drought}
\label{ch:te-heat-and-drought}
\hlchaptermodality{481}

\begin{chapteropening}
\textbf{By the end of this chapter:} \emph{I can say a heatwave, a mirage, a dust storm, to dry up and to crack open.}

Complete the last lesson of chapter 481: I can say a heatwave, a mirage, a dust storm, to dry up and to crack open.
\end{chapteropening}

\section[\texorpdfstring{\te{వడగాలి} (vaḍagāli)}{వడగాలి (vaḍagāli)}]{\te{వడగాలి} (vaḍagāli) — a heatwave; a hot wind}

\label{lesson:TE-C481-vadagali}



\begin{quote}
Before the new one: say the Telugu for to get bogged down, then the Telugu for to get stuck.
\end{quote}

\subsection*{You'll want to know: \te{వడగాలి}}
\textbf{\te{వడగాలి}} — \emph{vaḍagāli} — \textquotedblleft{}a heatwave; a hot wind\textquotedblright{}.

\begin{quote}
\te{ఎండాకాలంలో} \te{వడగాలి} \te{వీస్తుంది}.
\end{quote}

\begin{quote}
\emph{eṇḍākālaṁlō vaḍagāli vīstundi.}
\end{quote}

\begin{quote}
In the hot season a scorching wind blows.
\end{quote}

\te{వడ}, heat, + \te{గాలి}, wind. \te{వడదెబ్బ}, sunstroke, is the blow of the same heat.

\begin{cousinweb}
\textbf{\te{వడగాలి}} (\emph{vaḍagāli}) is native Dravidian: DEDR 5225 lists Telugu \emph{vaḍa}, \textquotedblleft{}heat\textquotedblright{}, and \emph{vaḍagāli}, \textquotedblleft{}the hot or land wind\textquotedblright{}.
\end{cousinweb}

\begin{grammarlens}[title={What you've built}]
One more word for getting around.
\end{grammarlens}

\subsection*{Your turn}
\begin{itemize}
\raggedright
  \item \emph{Say it:} \emph{vaḍagāli}
  \item \emph{Say it:} \emph{vaḍagāli}, once more
  \item \emph{Say it:} \emph{eṇḍākālaṁlō vaḍagāli vīstundi.} — In the hot season a scorching wind blows.
  \item \emph{Recall:} say the Telugu for to get bogged down, then the Telugu for to get stuck, then say \emph{vaḍagāli} again
\end{itemize}

\begin{tcolorbox}[breakable,colback=teal!4,colframe=teal!35!black,title={Before you move on}]
What does \te{వడగాలి} mean? (\textquotedblleft{}A heatwave; a hot wind\textquotedblright{}.) Say it once more.
\end{tcolorbox}
\section[\texorpdfstring{\te{ఎండమావి} (eṇḍamāvi)}{ఎండమావి (eṇḍamāvi)}]{\te{ఎండమావి} (eṇḍamāvi) — a mirage}

\label{lesson:TE-C481-endamavi}



\begin{quote}
Before the new one: say the Telugu for to get stuck, then the Telugu for a heatwave.
\end{quote}

\subsection*{You'll want to know: \te{ఎండమావి}}
\textbf{\te{ఎండమావి}} — \emph{eṇḍamāvi} — \textquotedblleft{}a mirage\textquotedblright{}.

\begin{quote}
\te{ఎడారిలో} \te{దూరంగా} \te{ఎండమావి} \te{కనిపించింది}.
\end{quote}

\begin{quote}
\emph{eḍārilō dūraṅgā eṇḍamāvi kanipiñcindi.}
\end{quote}

\begin{quote}
A mirage appeared far off in the desert.
\end{quote}

\te{ఎండ}, sunshine, + \te{మావి}. Something hoped for that is never reached is also an \te{ఎండమావి}.

\begin{cousinweb}
\textbf{\te{ఎండమావి}} (\emph{eṇḍamāvi}) is a Telugu compound: \te{ఎండ}, \textquotedblleft{}sunshine, the heat of the sun\textquotedblright{}, native Dravidian (DEDR 869), and -\te{మావి}, whose origin is not settled.
\end{cousinweb}

\begin{grammarlens}[title={What you've built}]
One more word for getting around.
\end{grammarlens}

\subsection*{Your turn}
\begin{itemize}
\raggedright
  \item \emph{Say it:} \emph{eṇḍamāvi}
  \item \emph{Say it:} \emph{eṇḍamāvi}, once more
  \item \emph{Say it:} \emph{eḍārilō dūraṅgā eṇḍamāvi kanipiñcindi.} — A mirage appeared far off in the desert.
  \item \emph{Recall:} say the Telugu for to get stuck, then the Telugu for a heatwave, then say \emph{eṇḍamāvi} again
\end{itemize}

\begin{tcolorbox}[breakable,colback=teal!4,colframe=teal!35!black,title={Before you move on}]
What does \te{ఎండమావి} mean? (\textquotedblleft{}A mirage\textquotedblright{}.) Say it once more.
\end{tcolorbox}
\section[\texorpdfstring{\te{గాలిదుమారం} (gālidumāraṁ)}{గాలిదుమారం (gālidumāraṁ)}]{\te{గాలిదుమారం} (gālidumāraṁ) — a dust storm}

\label{lesson:TE-C481-galidumaram}



\begin{quote}
Before the new one: say the Telugu for a heatwave, then the Telugu for a mirage.
\end{quote}

\subsection*{You'll want to know: \te{గాలిదుమారం}}
\textbf{\te{గాలిదుమారం}} — \emph{gālidumāraṁ} — \textquotedblleft{}a dust storm\textquotedblright{}.

\begin{quote}
\te{గాలిదుమారం} \te{వల్ల} \te{ఏమీ} \te{కనిపించలేదు}.
\end{quote}

\begin{quote}
\emph{gālidumāraṁ valla ēmī kanipiñcalēdu.}
\end{quote}

\begin{quote}
Because of the dust storm nothing could be seen.
\end{quote}

\te{గాలి}, wind, + \te{దుమారం}, a cloud of dust.

\begin{cousinweb}
\textbf{\te{గాలిదుమారం}} (\emph{gālidumāraṁ}) is a Telugu compound: \te{గాలి}, \textquotedblleft{}wind\textquotedblright{}, native Dravidian (DEDR 1499), and \te{దుమారం}, \textquotedblleft{}a whirl of dust\textquotedblright{}, whose origin is not settled.
\end{cousinweb}

\begin{grammarlens}[title={What you've built}]
One more word for getting around.
\end{grammarlens}

\subsection*{Your turn}
\begin{itemize}
\raggedright
  \item \emph{Say it:} \emph{gālidumāraṁ}
  \item \emph{Say it:} \emph{gālidumāraṁ}, once more
  \item \emph{Say it:} \emph{gālidumāraṁ valla ēmī kanipiñcalēdu.} — Because of the dust storm nothing could be seen.
  \item \emph{Recall:} say the Telugu for a heatwave, then the Telugu for a mirage, then say \emph{gālidumāraṁ} again
\end{itemize}

\begin{tcolorbox}[breakable,colback=teal!4,colframe=teal!35!black,title={Before you move on}]
What does \te{గాలిదుమారం} mean? (\textquotedblleft{}A dust storm\textquotedblright{}.) Say it once more.
\end{tcolorbox}
\section[\texorpdfstring{\te{ఎండిపోవు} (eṇḍipōvu)}{ఎండిపోవు (eṇḍipōvu)}]{\te{ఎండిపోవు} (eṇḍipōvu) — to dry up; to wither}

\label{lesson:TE-C481-endipovu}



\begin{quote}
Before the new one: say the Telugu for a mirage, then the Telugu for a dust storm.
\end{quote}

\subsection*{You'll want to know: \te{ఎండిపోవు}}
\textbf{\te{ఎండిపోవు}} — \emph{eṇḍipōvu} — \textquotedblleft{}to dry up; to wither\textquotedblright{}.

\begin{quote}
\te{కరువులో} \te{చెరువులన్నీ} \te{ఎండిపోయాయి}.
\end{quote}

\begin{quote}
\emph{karuvulō ceruvulannī eṇḍipōyāyi.}
\end{quote}

\begin{quote}
In the drought all the ponds dried up.
\end{quote}

Past (it …): \te{ఎండిపోయింది}, \emph{eṇḍipōyindi}.

\te{ఎండు}, to dry in the sun, + \te{పోవు}. Said of ponds, plants and throats. Its subject is a thing, so the past shown is the third person: \te{ఎండిపోయింది}, it dried up.

\begin{cousinweb}
\textbf{\te{ఎండిపోవు}} (\emph{eṇḍipōvu}) is native Dravidian: \te{ఎండు}, \textquotedblleft{}to be heated by the sun, to dry up\textquotedblright{} (DEDR 869), with \te{పోవు}, \textquotedblleft{}to go\textquotedblright{} (DEDR 4572).
\end{cousinweb}

\begin{grammarlens}[title={What you've built}]
One more word for getting around.
\end{grammarlens}

\subsection*{Your turn}
\begin{itemize}
\raggedright
  \item \emph{Say it:} \emph{eṇḍipōvu}
  \item \emph{Say it:} \emph{eṇḍipōvu}, once more
  \item \emph{Say it:} \emph{karuvulō ceruvulannī eṇḍipōyāyi.} — In the drought all the ponds dried up.
  \item \emph{Recall:} say the Telugu for a mirage, then the Telugu for a dust storm, then say \emph{eṇḍipōvu} again
\end{itemize}

\begin{tcolorbox}[breakable,colback=teal!4,colframe=teal!35!black,title={Before you move on}]
What does \te{ఎండిపోవు} mean? (\textquotedblleft{}To dry up; to wither\textquotedblright{}.) Say it once more.
\end{tcolorbox}
\section[\texorpdfstring{\te{బీటలు} \te{వారు} (bīṭalu vāru)}{బీటలు వారు (bīṭalu vāru)}]{\te{బీటలు} \te{వారు} (bīṭalu vāru) — to crack open}

\label{lesson:TE-C481-bitalu-varu}



\begin{quote}
Before the new one: say the Telugu for a dust storm, then the Telugu for to dry up.
\end{quote}

\subsection*{You'll want to know: \te{బీటలు} \te{వారు}}
\textbf{\te{బీటలు} \te{వారు}} — \emph{bīṭalu vāru} — \textquotedblleft{}to crack open\textquotedblright{}.

\begin{quote}
\te{ఎండలకు} \te{పొలాలు} \te{బీటలు} \te{వారాయి}.
\end{quote}

\begin{quote}
\emph{eṇḍalaku polālu bīṭalu vārāyi.}
\end{quote}

\begin{quote}
In the heat the fields cracked open.
\end{quote}

Past (it …): \te{బీటలు} \te{వారింది}, \emph{bīṭalu vārindi}.

\te{బీటలు}, cracks, + \te{వారు}. Said of dry earth and of old walls. Its subject is a thing, so the past shown is the third person: \te{బీటలు} \te{వారింది}, it cracked open.

\begin{cousinweb}
\textbf{\te{బీటలు} \te{వారు}} (\emph{bīṭalu vāru}) is native Dravidian: DEDR 5473 lists Telugu \emph{bīṭa}, \textquotedblleft{}a crack, a cleft\textquotedblright{}, and the phrase \emph{bīṭaluvāru}, \textquotedblleft{}to crack\textquotedblright{}, with Tamil \emph{veṭi}, \textquotedblleft{}to crack, to split\textquotedblright{}.
\end{cousinweb}

\begin{grammarlens}[title={What you've built}]
One more word for getting around.
\end{grammarlens}

\subsection*{Your turn}
\begin{itemize}
\raggedright
  \item \emph{Say it:} \emph{bīṭalu vāru}
  \item \emph{Say it:} \emph{bīṭalu vāru}, once more
  \item \emph{Say it:} \emph{eṇḍalaku polālu bīṭalu vārāyi.} — In the heat the fields cracked open.
  \item \emph{Recall:} say the Telugu for a dust storm, then the Telugu for to dry up, then say \emph{bīṭalu vāru} again
\end{itemize}

\begin{tcolorbox}[breakable,colback=teal!4,colframe=teal!35!black,title={Before you move on}]
What does \te{బీటలు} \te{వారు} mean? (\textquotedblleft{}To crack open\textquotedblright{}.) Say it once more.
\end{tcolorbox}
